\subsection{Design rationale: code-grounding over end-to-end control}
\label{sec:harness-rationale}

A vision-language-action (VLA) model collapses perception, reasoning,
and control into one policy, trained or fine-tuned end-to-end on
embodiment-specific demonstration data, which is the bottleneck: paired
vision-action trajectories remain far scarcer and more expensive to
collect than the text, code, and image-text corpora a general-purpose
LLM or VLM is pretrained on. A decoupled Sense-Plan-Act architecture
sidesteps this. If planning is text (or text conditioned on an image),
and the resulting plan is a sequence of parameterized skills rather than
raw joint torques, then whatever manipulation-relevant knowledge the
model already has, spatial reasoning, arithmetic, common-sense physics,
ordering and precondition logic, is available zero-shot, through code,
without a single task-specific demonstration, the same argument prior
code-generation-for-robotics work makes~\cite{liang2023codeaspolicies,
huang2023voxposer, huang2022innermonologue, singh2023progprompt}: an LLM
that has never seen a robot trajectory has still seen a great deal of
code.

The cost is that a decoupled architecture asks the model to compute
physical quantities, a release height, a grasp width, a target
position, that it did not learn from a source grounded in the real
world at all; \S\ref{sec:taxonomy} is what happens when that
computation is trusted uncritically. Our answer is not to abandon the
decoupled architecture but to constrain it: let the model choose
\emph{what} to do and \emph{where}, symbolically, and let a
deterministic layer compute \emph{how much}, from a sensor measurement
or a paired action already in the plan, whenever the plan needs a
number the model cannot be trusted to produce on its own.

\subsection{System architecture}
\label{sec:harness-architecture}

\begin{figure*}[t]
\centering
\fbox{\begin{minipage}[c][5.5cm]{0.92\textwidth}
\centering
System architecture diagram: task interface/GUI, the three planner
nodes, the plan manager, the skill executor, and the code-grounding
layer, with the debug-artifact and configuration paths annotated.
\\[4pt]
\footnotesize (placeholder; final figure to be supplied)
\end{minipage}}
\caption{RoboReason harness architecture: the Sense-Plan-Act pipeline of
\S\ref{sec:harness-architecture}, with the code-grounding layer overriding
model-computed geometry before execution.}
\label{fig:architecture}
\end{figure*}

The harness is a ROS2 stack organized around the same
Sense-Plan-Act pipeline as \S\ref{sec:related-work}, shown in
Figure~\ref{fig:architecture}. A task interface (a
terminal client or a web GUI) forwards a natural-language command to one
of three planner nodes: a static-scene LLM planner, a camera-grounded VLM
planner, or a hybrid that runs a VLM scene-grounding call first and hands
the resulting scene to the same LLM planner used in static-scene mode.
Whichever planner produced it, the plan is validated and sequenced by a
plan-manager node, which drives a skill-executor node (a physics-free
mock for fast iteration, or the real UR5cb driver) through a small,
fixed set of parameterized skills: \code{approach}, \code{pick},
\code{release}, \code{move\_home}, and \code{wait}. Every planning call
supports six interchangeable reasoning strategies, from a single-step
reactive baseline to self-consistency voting and tree search, selectable
without touching the pipeline itself.

Two design choices let this pipeline scale into a research tool rather
than a one-off demo. First, every runtime parameter, model choice,
reasoning strategy, grounding modality, timeouts, robot IP, is a field
on one centralized settings object, overridable by an environment
variable with no code change, which is what lets
\S\ref{sec:evaluation}'s conditions be defined and reproduced as
configuration rather than as code branches. Second, every planning call
records a per-run debug artifact, the command, active configuration, raw
model response, execution log, and, in VLM modes, the camera frame the
plan was grounded in, which is what made diagnosing
\S\ref{sec:taxonomy}'s failures possible: each bug was traced from a
hardware failure back to a specific malformed field in a specific logged
response, not inferred after the fact. A web GUI supervises the whole
stack, including the physical UR driver and camera service, both of
which fail intermittently enough to need their own restart logic, as
monitored child processes, so a trial can be launched, watched, and
annotated without a terminal per node.
